\documentclass[12pt]{article}
\usepackage[utf8]{inputenc}
\usepackage[T1]{fontenc}
\usepackage[margin=1in]{geometry}
\usepackage{setspace}
\usepackage{xcolor}
\usepackage[colorlinks=true, urlcolor=blue!60!black, linkcolor=blue!60!black]{hyperref}
\setlength{\parindent}{0em}
\setlength{\parskip}{0.6em}
\renewcommand{\baselinestretch}{1.0}
\begin{document}
\begin{center}
{\LARGE \textit{United States v. Belmont}, 301 U.S. 324 (1937)\par}
\vspace{0.5em}
{\large Case Brief --- Two-Page Edition\par}
\end{center}

\section{Memory Jogger}
A sole executive agreement made incident to recognizing the Soviet Union needed no Senate consent and overrode New York's public policy (\href{https://supreme.justia.com/cases/federal/us/301/324/}{Justia}).

\section{Detailed Case Facts}
After the Bolshevik Revolution, the Soviet Government nationalized a Russian corporation and its New York bank deposit with banker August Belmont. When President Roosevelt recognized the Soviet Union in 1933, the Litvinov Assignment, never submitted to the Senate, assigned Soviet claims against Americans to the United States.

\section{Procedural History}
The lower courts dismissed the United States' suit against Belmont's executors, holding that enforcing the confiscation would violate New York public policy.

\section{Judicial Votes}
Reversed 9--0 in the judgment. Sutherland wrote for six Justices; Stone, joined by Brandeis and Cardozo, concurred in the judgment.

\section{Holding}
Recognition validated the Soviet Government's acts; the President could make the assignment without the Senate; and a valid sole executive agreement overrides conflicting state law and policy.

\section{Analysis of Opinions}
\textbf{Majority:} the recognition power includes settling claims that obstruct relations; compacts such as protocols and claims assignments need no Senate consent; in foreign affairs ``state lines disappear''; and the Constitution does not protect a foreign corporation abroad. \textbf{Stone:} New York had no policy letting a mere stakeholder refuse payment, and the agreement did not purport to override state law.

\section{Examples: Future Applications}
\textbf{Same-Side:} when a claims settlement accompanies recognition of a new government, a contrary state statute yields. \textbf{Opposite-Side:} an agreement directing state criminal courts, without legislation, fails (\textit{Medell\'in}). \textbf{Fence-Sitter:} ending state tort suits without a compensation fund or recognition.

\section{Critique}
Letting a sole executive agreement preempt state law bypasses the Treaty Clause and Article VI's text; Stone's narrower approach, later echoed in \textit{Medell\'in}, avoided that maximalism.

\section{Key Quotations}
\begin{quote}
``Plainly, the external powers of the United States are to be exercised without regard to state laws or policies.''
\end{quote}
\noindent --- Justice Sutherland, for the Court, 301 U.S. at 331.

\end{document}
